\section{渐近区间与 Bayesian 可信区间}
\label{sec:11-Mathematical-Statistics/04-Interval-Estimation/03}

上一节构造置信区间的方法依赖于一个精确定量——枢轴量，其分布完全已知，不依赖任何未知参数。\(\bar X\) 减 \(\mu\) 除以 \(S/\sqrt{n}\) 恰好是 \(t\) 分布，这个巧合是 Gauss 模型的礼物。多数模型没有这份运气：似然函数没有简洁的充分统计量，或者估计量的精确分布难以追踪。

难题很清楚：我们手上有一个估计量 \(\hat\theta\)，知道它在某种渐近意义下``逼近''正态分布，但要对有限样本给出一个区间声明。大样本理论提供了多条路，每一条都在不同的地方做了近似，在不同的地方可能翻车。

\subsection{最直接的路：Wald 区间}

若估计量满足
\[
\frac{\hat\theta-\theta}{\widehat{\SE}(\hat\theta)}\overset d\to N(0,1),
\]
以正态分位数构造对称区间：
\[
\hat\theta\pm z_{1-\alpha/2}\,\widehat{\SE}(\hat\theta).
\]
构造极其简单：一个点估计，加减若干倍标准误。这使它成为几乎所有统计软件默认输出的区间。

但它的简单是有代价的。标准误 \(\widehat{\SE}\) 本身也是估计量，把未知 \(\theta\) 替换为 \(\hat\theta\)（plug-in）引入了第二层近似。更重要的是，Wald 区间把抽样分布强制当作关于 \(\hat\theta\) 对称的正态曲线——在参数空间内部、大样本下，这个近似常常还行；到了参数边界附近，或者样本不够大时，真实分布可能严重偏斜，对称区间的覆盖率就会出现系统偏差。

第一节的 Bernoulli 例子正是这个失败的生动写照：\(\theta\) 接近 0 或 1 时，Wald 区间的边界可能超出 \([0,1]\)，覆盖率远低于名义水平。Wald 区间方便，但方便的前提条件被违反时，它不会给你任何警告。

\subsection{Delta method：把不确定性传给变换后的参数}

对估计量做非线性变换之后，不能简单地把点估计变换一下、标准误原封不动地套上去。设
\[
\sqrt n(\hat\theta-\theta)\overset d\to N(0,V),
\]
而 \(g\) 在 \(\theta\) 处可微。Taylor 展开到一阶：
\[
g(\hat\theta)-g(\theta)\approx g'(\theta)(\hat\theta-\theta),
\]
由此
\[
\sqrt n\bigl(g(\hat\theta)-g(\theta)\bigr)\overset d\to N\bigl(0,[g'(\theta)]^2V\bigr).
\]
方差被 \(g'(\theta)\) 的平方缩放了一步。实际操作中，渐近方差里的未知 \(\theta\) 替换为 \(\hat\theta\)（plug-in），于是 \(g(\hat\theta)\) 的标准误近似为
\[
|g'(\hat\theta)|\,\widehat{\SE}(\hat\theta).
\]

Delta method 的一个常见妙用是在 log 尺度上构造区间。对于只能取正值的参数（如方差、率参数），先在 log 尺度上构造一个对称的正态近似区间，再指数变换回原尺度。这样得到的区间自动保持正性，而且能更好地适应乘法性质的不确定性——因为 log 变换把乘法模型变成了加法模型，Wald 在加法尺度上的正态近似更合理。

若 \(g'(\theta)=0\)，一阶 delta method 退化，主导项变成二阶导数，需要二阶展开或其他渐近理论。函数在参数点不可微时也是如此——一阶 Taylor 的前提从根上就不成立。

\subsection{Profile likelihood 区间：让似然函数自己说话}

当模型含有不感兴趣的参数（nuisance parameter）时，Wald 区间需要对整个参数向量做联合渐近正态近似，太粗糙。Profile likelihood 走另一条路：对感兴趣的参数 \(\psi\) 的每一个候选值，把其余参数在固定 \(\psi\) 的条件下最大化，得到 profile log-likelihood
\[
\ell_p(\psi)=\max_{\lambda}\ell(\psi,\lambda).
\]
直观上，你问``这个模型在 \(\psi\) 取某个具体值时，最多能有多大的似然''——把最好情况下的拟合度作为 \(\psi\) 的函数。

在正则条件（Wilks 定理，假设检验部分将系统讨论）下，似然比统计量
\[
2\bigl[\ell_p(\hat\psi)-\ell_p(\psi)\bigr]
\]
在大样本下近似 \(\chi^2_1\) 分布。按上一节的检验—区间对偶，收集所有未被拒绝的 \(\psi\) 值，就得到 profile likelihood 置信区间。

它的优势在于直接反映似然函数的形状：似然偏斜多厉害，区间就偏斜多厉害；参数换一种变换，区间自动跟着变换，不需要另行计算导数。代价是每一个候选 \(\psi\) 都需要做一次受约束优化，计算量比 Wald 区间的点估计加减标准误大得多。边界参数、混合模型和弱识别会改变极限分布，\(\chi^2_1\) 近似在这些场景不再可靠。

Bootstrap 区间提供了第三条路：通过重采样近似估计量的整个抽样分布，不依赖显式的渐近方差公式。Percentile、basic 与 BCa 等方法将在重采样部分专门讨论。Bootstrap 也不是免费午餐——厚尾、强依赖和非光滑统计量都可能让它失效。

\subsection{Bayesian 可信区间：换一种概率问题的答案}

前面所有区间都回答同一个问题：``如果真实参数是 \(\theta\)，数据落在某个范围内的概率是多少？''这是频率学派的问题。Bayesian 换了一个问题：``数据已经是眼前这组数了，参数落在某个范围内的概率是多少？''

给定先验分布和数据，后验区间 \(C(x)\) 满足
\[
P\{\Theta\in C(x)\mid x\}=1-\alpha.
\]
这里的概率来自参数的后验分布，数据 \(x\) 是固定的。于是，在特定模型和先验下，可以直接说：``参数落在该区间内的后验概率为 95\%''。这个表述正是非统计专业人士期待听到的那种话。

常用的可信区间构造方式有三种：

\begin{itemize}
\tightlist
\item
  等尾区间：后验分布的左右两侧各留下 \(\alpha/2\) 概率，区间取中间。实现简单，对单峰对称后验合理。
\item
  最高后验密度（HPD）区间：取后验密度最高的 \(1-\alpha\) 区域。在给定的覆盖率下做到区间最短，但对多峰后验可能返回不连通区域——挑出几个密度山峰，中间的低谷不要。如果此时强行报告单一连续区间，山谷里那些不太可能的参数值反而被装了进去。
\item
  单侧可信界：只关心参数不低于或不高于某个阈值时用。
\end{itemize}

三种选择可能给出形状不同的区间，却都标着相同的后验概率。报告可信区间时必须说明类型，否则读者无法判断区间的构造逻辑。

\subsection{可信度与覆盖率不能互换}

Bayesian 可信区间有后验概率解释，但它不对每个固定参数提供 95\% 的频率覆盖保证。反过来，频率置信区间提供重复抽样下的覆盖率，却不能给这次具体数据下的``参数概率''。

在正则大样本下，Bernstein--von Mises 定理常使两类区间在数值上接近——先验的影响被数据淹没，后验分布近似于以 MLE 为中心的正态分布，和频率渐近区间重合。但以下场景中差异可以很大：小样本（先验的权重不可忽略）、强先验（数据不足以压倒先验）、边界参数（渐近正态性不成立）、高维问题（Bernstein--von Mises 可能需要不现实的样本量）。

面对一个区间声明，应该能区分五种不同的概率陈述：

\begin{itemize}
\tightlist
\item
  概率是对数据重复而言（频率），还是对参数后验而言（Bayesian）？
\item
  覆盖是逐参数保证、对先验平均、还是仅在大样本下渐近成立？
\item
  先验有没有被数据反复挑选——即先验本身是否已经是看过数据后调整的结果？
\item
  区间是不是在模型选择之后才构造的——模型选择的不确定性有没有进入区间宽度？
\item
  参数化方式改变后，区间对应变换是否与参数的实际不确定性传播一致？
\end{itemize}

两类区间不需要说成互相排斥。很多实际问题中，同时报告频率置信区间和 Bayesian 可信区间并说明差异来源，比宣称其中一方是``正确''的更有信息量。但无论如何，不要把两种概率解释混在同一条区间声明里——``这个区间以 95\% 概率包含参数''既不是严格的频率陈述，也不是严格的 Bayesian 陈述，它只是把两边的好处各取一半拼成了一句统计上不成立的句子。
